\chapter{Psychiatry}

\medterm{Acute Stress Disorder}-- Characterized by the development of specific symptoms following exposure to a traumatic event. The same symptom clusters as PTSD are required: intrusion, negative mood, dissociation, avoidance, and arousal. Symptoms must last for at least three days and up to one month after the trauma. If symptoms persist beyond one month, the diagnosis changes to PTSD. Early intervention may prevent progression to PTSD.

\medterm{Addiction}-- A chronic, relapsing brain disorder characterized by compulsive substance use or behavior despite harmful consequences. It involves changes in brain reward, motivation, and memory circuits, leading to impaired control over the behavior. Treatment often combines medication, behavioral therapy, and social support.

\medterm{Adjustment Disorder}-- Emotional or behavioral symptoms in response to an identifiable stressor occurring within three months of the onset of the stressor. The symptoms are disproportionate to the severity of the stressor and cause significant distress or impairment. It resolves within six months of the termination of the stressor. Subtypes include with depressed mood, with anxiety, with mixed anxiety and depressed mood, with disturbance of conduct, with mixed disturbance of emotions, and unspecified.

\medterm{Affective Flattening}-- A significant reduction in the range and intensity of emotional expression. It is a negative symptom of schizophrenia characterized by reduced facial expression, monotonous voice, and decreased body language. It is distinct from blunted affect (reduced range) and constricted affect (narrowed range). It is less responsive to antipsychotics.

\medterm{Agitation}-- A state of excessive motor activity, emotional excitement, and restlessness. In psychiatric settings, it may be a symptom of mania, psychosis, delirium, substance intoxication, or withdrawal. Acute agitation is managed with de-escalation techniques first, then pharmacological management with benzodiazepines (lorazepam), antipsychotics (haloperidol, olanzapine IM), or combination (B52: diphenhydramine + haloperidol + lorazepam). Physical restraints are used as a last resort.

\medterm{Agoraphobia}-- Marked fear or anxiety about two or more of five situations: using public transportation, being in open spaces, being in enclosed places, standing in line or being in a crowd, and being outside the home alone. The individual fears or avoids these situations because escape might be difficult or help unavailable if panic-like symptoms develop. The fear is out of proportion, persists for six months, and causes significant impairment.

\medterm{Alcohol Use Disorder}-- A problematic pattern of alcohol use leading to significant impairment or distress. Severity is classified as mild (2--3 criteria), moderate (4--5 criteria), or severe (6 or more criteria). Criteria include drinking more or longer than intended, persistent desire to cut down, spending excessive time with alcohol, craving, failure to fulfill obligations, continued use despite social or interpersonal problems, and tolerance and withdrawal. Withdrawal can include seizures, delirium tremens, and autonomic hyperactivity. Severe alcohol withdrawal can progress to delirium tremens, a life-threatening condition with a mortality rate of 5-15 percent if untreated. Management of alcohol withdrawal involves the CIWA-Ar (Clinical Institute Withdrawal Assessment for Alcohol, revised) scale to guide benzodiazepine administration. First-line treatment is symptom-triggered benzodiazepine therapy (lorazepam, diazepam, or chlordiazepoxide). Supportive care includes intravenous fluids, thiamine supplementation to prevent Wernicke-Korsakoff syndrome (100 mg IV or IM daily), correction of electrolyte disturbances (hypokalaemia, hypomagnesaemia, hypophosphataemia), and multivitamins. Carbamazepine and gabapentin are alternative treatments for mild-to-moderate withdrawal. Severe withdrawal may require ICU admission and management of delirium tremens.

\medterm{Alcohol Withdrawal}-- Symptoms occurring after cessation or reduction of prolonged heavy drinking. Timeline includes: tremor and anxiety at 6--12 hours, perceptual disturbances at 12--24 hours, withdrawal seizures at 24--48 hours, and delirium tremens at 48--96 hours. Delirium tremens includes autonomic hyperactivity, tremor, hallucinations, and confusion with a mortality rate of 5--15\% if untreated. CIWA-Ar scale guides management. Benzodiazepines are first-line treatment.

\medterm{Alogia}-- Poverty of speech or poverty of content of speech, a negative symptom of schizophrenia. It manifests as brief, empty replies to questions, reduced spontaneous speech, and lack of elaboration. It reflects dysfunction in the mesocortical dopamine pathway and is less responsive to antipsychotic medications than positive symptoms.

\medterm{Anhedonia}-- The inability to experience pleasure from activities that were previously enjoyable. It is a core symptom of major depressive disorder and a negative symptom of schizophrenia. It is mediated by dysfunction in the brain's reward circuitry, particularly the mesolimbic dopamine system.

\medterm{Anorgasmia}-- Persistent or recurrent delay in, or absence of, orgasm following a normal sexual excitement phase. Causes: antidepressants (especially SSRIs, the most common iatrogenic cause), psychosocial factors (anxiety, relationship issues, past trauma), neurological disorders (spinal cord injury, multiple sclerosis, diabetic neuropathy), hormonal abnormalities (low testosterone, hypothyroidism), and pelvic surgery. Management: for SSRI-induced anorgasmia, strategies include dose reduction, drug holiday, switching to bupropion or mirtazapine, or augmentation with buspirone, cyproheptadine, or sildenafil. Psychosexual therapy and behavioural techniques are effective for psychogenic causes. Prognosis depends on etiology; medication-induced anorgasmia is usually reversible.

\medterm{Anticholinergic Burden}-- The cumulative effect of taking multiple medications with anticholinergic properties. It can cause dry mouth, constipation, urinary retention, blurred vision, confusion, and delirium, especially in elderly patients. Many psychiatric medications (TCAs, first-generation antipsychotics, some antihistamines) have anticholinergic effects.

\medterm{Anticonvulsant Mood Stabilizers}-- Antiepileptic drugs used as mood stabilizers in bipolar disorder. Valproate (divalproex) is effective for acute mania and rapid cycling. Carbamazepine is effective for acute mania but has numerous drug interactions via CYP3A4 induction. Lamotrigine is effective for bipolar depression prevention but not acute mania. All require monitoring for side effects and drug levels (except lamotrigine, which is dose-titrated).

\medterm{Antisocial Personality Disorder}-- A pattern of disregard for and violation of the rights of others occurring since age 15 years. At least three of seven criteria: failure to conform to social norms, deceitfulness, impulsivity, irritability and aggressiveness, reckless disregard for safety, consistent irresponsibility, and lack of remorse. Diagnosis requires evidence of conduct disorder before age 15. Associated with psychopathy, which includes affective and interpersonal features beyond the DSM criteria.

\medterm{Anxiety}-- A normal emotional response to perceived danger, but anxiety disorders involve excessive, persistent worry that interferes with daily life. Symptoms can include restlessness, rapid heart rate, muscle tension, and avoidance behaviors. Common types include generalized anxiety disorder, panic disorder, social anxiety, and phobias, all of which are treatable with therapy and medication.

\medterm{Anxiety Disorder}-- A group of psychiatric conditions characterized by excessive fear, anxiety, and related behavioral disturbances. Includes generalized anxiety disorder (GAD), panic disorder, social anxiety disorder, specific phobias, and agoraphobia. GAD: persistent, excessive worry about multiple domains for \(\ge\)6 months, with restlessness, fatigue, difficulty concentrating, irritability, muscle tension, and sleep disturbance. Treatment: SSRIs/SNRIs (sertraline, escitalopram, venlafaxine), buspirone, benzodiazepines for short-term use, and cognitive-behavioral therapy (CBT).

\medterm{Applied Behavior Analysis}-- A behavioral intervention based on learning theory principles that uses systematic reinforcement to shape desired behaviors and reduce problematic behaviors. It is the primary evidence-based intervention for autism spectrum disorder, particularly in young children. Core components include positive reinforcement, prompting, fading, chaining, and generalization. It requires intensive, structured sessions.

\medterm{Attention-Deficit/Hyperactivity Disorder}-- A neurodevelopmental disorder characterized by persistent patterns of inattention, hyperactivity, and impulsivity that interfere with functioning or development. Three presentations: predominantly inattentive, predominantly hyperactive-impulsive, and combined. Inattention symptoms include difficulty sustaining attention, poor organization, forgetfulness, and distractibility. Hyperactivity-impulsive symptoms include fidgeting, excessive talking, difficulty waiting turns, and interrupting. Must have symptoms present before age 12. Diagnosis is clinical, using DSM-5 criteria, supported by rating scales (Conners Rating Scales, SNAP-IV, Vanderbilt Assessment Scale), collateral history from parents and teachers, and neuropsychological testing to evaluate attention and executive function. Management includes pharmacological treatment with psychostimulants (methylphenidate, dexamphetamine, lisdexamfetamine) as first-line, non-stimulant options (atomoxetine, guanfacine, clonidine), and behavioural interventions (parent training, behavioural classroom management, organisational skills training). Prognosis with treatment is good; untreated ADHD is associated with academic underachievement, social difficulties, substance misuse, and occupational impairment in adulthood.

\medterm{Attention-Deficit/Hyperactivity Disorder (ADHD)}-- A neurodevelopmental condition marked by persistent inattention, hyperactivity, and impulsivity that interferes with daily functioning. In children, symptoms often include difficulty staying seated, frequent interrupting, and trouble sustaining attention. Adults may experience restlessness, poor time management, and difficulty concentrating, though hyperactivity often diminishes with age.

\medterm{Atypical Features}-- A specifier for depressive episodes characterized by mood reactivity, significant weight gain or increased appetite, hypersomnia, leaden paralysis, and a long-standing pattern of interpersonal rejection sensitivity. Atypical depression often presents in younger patients and may respond preferentially to MAOIs or SSRIs rather than TCAs.

\medterm{Autism Spectrum Disorder}-- A neurodevelopmental disorder characterized by persistent deficits in social communication and social interaction across multiple contexts, and restricted, repetitive patterns of behavior, interests, or activities. Social deficits include difficulty with social-emotional reciprocity, nonverbal communicative behaviors, and developing/maintaining relationships. Restricted behaviors include stereotyped movements, insistence on sameness, and restricted interests. Severity is specified both domains. Diagnosis is based on DSM-5 criteria. Management includes early intensive behavioural intervention (applied behaviour analysis), speech and language therapy, social skills training, occupational therapy, and pharmacological treatment of associated conditions such as SSRIs for anxiety and antipsychotics for irritability. Prognosis is variable, with early intervention improving outcomes.

\medterm{Avoidant Personality Disorder}-- A pattern of social inhibition, feelings of inadequacy, and hypersensitivity to negative evaluation. At least four of seven criteria: avoids occupational activities involving interpersonal contact, unwilling to get involved unless certain of being liked, shows restraint within intimate relationships due to fear of shame, preoccupied with being criticized or rejected in social situations, inhibited in new interpersonal situations due to feelings of inadequacy, views self as socially inept, and is unwilling to get involved unless certain of being liked. Diagnosis is clinical using DSM-5 criteria. Treatment includes psychodynamic therapy and CBT focusing on social skills training and graded exposure to feared social situations. Pharmacotherapy with SSRIs may help with comorbid social anxiety. Prognosis is guarded, as the disorder is typically lifelong, though symptoms may improve with treatment.

\medterm{Avolition}-- A severe lack of motivation to initiate and perform self-directed, purposeful activities. It is a negative symptom of schizophrenia and is seen in severe depression. It manifests as neglect of personal hygiene, inability to complete tasks, and social withdrawal. It is mediated by dysfunction in the prefrontal cortex and dopamine reward pathways.

\medterm{Benzodiazepine Use Disorder}-- A problematic pattern of benzodiazepine use leading to significant impairment. Benzodiazepines include alprazolam, clonazepam, diazepam, and lorazepam. Dependence can develop rapidly with regular use. Withdrawal can be life-threatening, presenting with anxiety, insomnia, tremor, seizures, and psychosis. Cross-tolerance with alcohol exists. Abrupt discontinuation should be avoided; gradual taper is essential.

\medterm{Benzodiazepine Withdrawal}-- Symptoms from cessation or reduction of prolonged benzodiazepine use. Early symptoms include anxiety, insomnia, tremor, and diaphoresis. Severe withdrawal can include seizures, psychosis, and delirium. Withdrawal severity depends on the half-life of the benzodiazepine, duration of use, and dose. Abrupt discontinuation after chronic use can be fatal. A slow taper, often switching to a long-acting benzodiazepine, is the standard approach.

\medterm{Bipolar Disorder}-- A mood disorder characterized by episodes of mania/hypomania and depression. Bipolar I: at least one manic episode (elevated mood, grandiosity, decreased need for sleep, pressured speech, increased goal-directed activity, risky behavior) lasting \(\ge\)1 week or requiring hospitalization. Bipolar II: hypomanic episodes and major depressive episodes without full mania. Treatment: mood stabilizers (lithium, valproate, lamotrigine), atypical antipsychotics, and psychotherapy. Antidepressants used cautiously due to risk of mood switching.

\medterm{Bipolar I Disorder}-- A mood disorder defined by at least one manic episode lasting at least seven days or requiring hospitalization. A manic episode is a distinct period of abnormally elevated, expansive, or irritable mood and increased goal-directed activity or energy. Associated symptoms include inflated self-esteem, decreased need for sleep, pressured speech, flight of ideas, distractibility, increased activity, and risky behavior. Major depressive episodes may occur but are not required for diagnosis.

\medterm{Bipolar II Disorder}-- A mood disorder characterized by at least one hypomanic episode and at least one major depressive episode, with no history of a full manic episode. The depressive episodes in bipolar II tend to be more frequent and longer-lasting than in bipolar I, leading to greater cumulative disability. Patients may be misdiagnosed with unipolar depression if the hypomanic episodes are not identified.

\medterm{Body Dysmorphic Disorder}-- Preoccupation with one or more perceived defects or flaws in physical appearance that are not observable or appear slight to others. The individual performs repetitive behaviors in response to appearance concerns such as mirror checking, excessive grooming, skin picking, reassurance seeking, and comparing. The preoccupation causes clinically significant distress or impairment. Insight is often poor, and patients may seek cosmetic procedures that do not resolve the concern.

\medterm{Borderline Personality Disorder}-- A pattern of instability in interpersonal relationships, self-image, and affects, with marked impulsivity. At least five of nine criteria: frantic efforts to avoid abandonment, unstable and intense relationships, identity disturbance, impulsivity, recurrent suicidal behavior, affective instability, chronic feelings of emptiness, inappropriate intense anger, and transient paranoid ideation or dissociation. Strongly associated with childhood trauma, especially sexual abuse. DBT is the evidence-based treatment of choice. Other treatments include mentalisation-based therapy (MBT), transference-focused psychotherapy (TFP), and general psychiatric management (GPM). Pharmacotherapy targets specific symptoms: SSRIs for mood instability, low-dose antipsychotics for cognitive-perceptual symptoms, and mood stabilisers for affective dysregulation. Prognosis is guarded; symptoms typically improve with age and treatment, but functional impairment often persists.

\medterm{Brief Psychotic Disorder}-- A disturbance characterized by the sudden onset of at least one psychotic symptom: delusions, hallucinations, disorganized speech, or grossly disorganized or catatonic behavior. The disturbance lasts at least one day but less than one month, after which the individual fully returns to premorbid functioning. It may occur in response to a very stressful event and is diagnosed only when the disturbance is not better explained by another psychotic disorder, mood disorder, or substance use.

\medterm{Caffeine Withdrawal}-- A syndrome occurring after abrupt cessation or reduction of caffeine use. Symptoms include headache, fatigue, depressed mood, difficulty concentrating, and flu-like symptoms. Onset is typically 12--24 hours after last caffeine intake, peak at 20--51 hours, and duration of 2--9 days. It is one of the most commonly experienced withdrawal syndromes.

\medterm{Cannabis Use Disorder}-- A problematic pattern of cannabis use leading to significant impairment. Criteria include tolerance and withdrawal (irritability, anger, decreased appetite, restlessness, anxiety, sleep disturbance). Cannabis intoxication includes euphoria, impaired judgment, and time distortion. High-potency concentrates (dabbing) increase the risk of dependence. Cannabinoid hyperemesis syndrome presents with cyclic vomiting in heavy users.

\medterm{Catatonia}-- A neuropsychiatric syndrome characterized by motor abnormalities including immobility, mutism, posturing, waxy flexibility, stereotypy, mannerisms, grimacing, echolalia, and echopraxia. It can occur in mood disorders, psychotic disorders, medical conditions, or as catatonia due to another medical condition. Lorazepam challenge and ECT are first-line treatments. Immobility and autonomic instability suggest malignant catatonia.

\medterm{Claustrophobia}-- An anxiety disorder characterized by an intense, irrational fear of enclosed or confined spaces (e.g., lifts, tunnels, MRI scanners, windowless rooms). Individuals experience marked anxiety, panic attacks, and avoidance behaviour. Prevalence: approximately 5-10\% of the general population. Claustrophobia is classified as a specific phobia (situational type) in DSM-5. The fear is out of proportion to the actual danger and causes significant distress or functional impairment. Pathophysiology involves conditioned fear responses, catastrophic misinterpretation of bodily sensations, and possibly a heightened sensitivity to suffocation cues. Management: CBT with graded in vivo exposure is the most effective treatment; virtual reality exposure therapy is an emerging alternative. Pharmacotherapy (SSRIs, benzodiazepines) may be used adjunctively. Prognosis is excellent with treatment; most patients achieve significant symptom reduction.

\medterm{Cognitive Behavioral Therapy}-- A structured, time-limited psychotherapy that identifies and modifies maladaptive thought patterns (cognitive distortions) and behaviors. It is based on the cognitive model that thoughts, feelings, and behaviors are interconnected. CBT uses techniques including cognitive restructuring, behavioral activation, exposure, and skills training. It is the most evidence-based psychotherapy for depression, anxiety disorders, OCD, PTSD, and eating disorders.

\medterm{Cognitive Behavioral Therapy for Insomnia (CBT-I)}-- Cognitive behavioral therapy for insomnia is a structured, evidence-based program that targets the thoughts and behaviors perpetuating poor sleep. Components include sleep restriction, stimulus control, cognitive restructuring, and sleep hygiene education. CBT-I is recommended as first-line treatment for chronic insomnia and yields durable improvements without the risks of sleep medications.

\medterm{Cognitive Distortions}-- Systematic errors in thinking that perpetuate negative emotional states and maladaptive behaviors. Common distortions include all-or-nothing thinking, catastrophizing, mind reading, overgeneralization, emotional reasoning, and should statements. CBT focuses on identifying and modifying these patterns.

\medterm{Comorbidity}-- The co-occurrence of two or more psychiatric disorders in a single patient. Psychiatric comorbidity is the rule rather than the exception. For example, depression commonly co-occurs with anxiety disorders, substance use disorders, and personality disorders. Comorbidity complicates treatment and worsens prognosis.

\medterm{Compulsions}-- Repetitive behaviors or mental acts that the individual feels driven to perform in response to an obsession or according to rigid rules. They are aimed at preventing or reducing anxiety or preventing some dreaded event, but are not connected in a realistic way with what they are designed to neutralize. Common compulsions include washing, checking, counting, ordering, arranging, and mental rituals. Time-consuming compulsions (more than one hour daily) are associated with greater severity.

\medterm{Conduct Disorder}-- A childhood/adolescent disorder characterized by a persistent pattern of violating others' basic rights or major age-appropriate norms. Behaviors: aggression to people/animals, destruction of property, deceitfulness/theft, and serious rule violations. Diagnosis requires \(\ge\)3 criteria in the past 12 months with functional impairment. Risk factors: abuse, neglect, parental substance use, and genetic predisposition. Treatment: parent management training, multisystemic therapy, and CBT; limited pharmacotherapy for comorbid conditions.

\medterm{Consultation-Liaison Psychiatry}-- The subspecialty of psychiatry focused on the interface between psychiatry and general medical illness. Psychiatrists provide consultation on medical-surgical wards for patients with psychiatric symptoms arising from or coexisting with medical conditions. Common issues include delirium, depression in medically ill patients, factitious disorder, adjustment reactions, and capacity assessments. It requires expertise in both psychiatric and medical conditions.

\medterm{Conversion Disorder}-- A condition in which patients present with neurological symptoms (weakness, paralysis, non-epileptic seizures, sensory loss, gait abnormalities, dystonia, speech difficulties) that cannot be explained by organic disease. Previously known as hysteria. The symptoms are not intentionally produced (unlike malingering or factitious disorder). Associated with psychological stressors and trauma. Most common in young women. Diagnosis: positive signs on neurological exam (Hoover sign for weakness, variability in dystonia), inconsistency with known neurological disease. Treatment: psychoeducation, physical therapy, CBT, and treatment of comorbid psychiatric conditions. Prognosis: variable; acute onset has better outcomes.

\medterm{Countertransference}-- The therapist's emotional reactions to the patient, which may be influenced by the therapist's own unresolved conflicts or may be a response to the patient's transference. Awareness of countertransference is essential for effective therapy and can provide valuable clinical information about the patient's interpersonal style.

\medterm{Cross-Tolerance}-- Tolerance to one substance that extends to another substance with a similar mechanism of action. For example, tolerance to one benzodiazepine extends to all benzodiazepines, and tolerance to alcohol extends partially to benzodiazepines due to shared GABA-A receptor mechanisms. Cross-tolerance is clinically relevant when switching between substances of the same class.

\medterm{Cultural Psychiatry}-- The study and clinical practice of how cultural factors influence mental health, illness expression, diagnosis, and treatment. Cultural concepts of distress, idioms of distress, and cultural syndromes must be considered. The DSM-5 includes a Cultural Formulation Interview. Culture influences symptom presentation, help-seeking behavior, treatment adherence, and the therapeutic relationship.

\medterm{Cyclothymic Disorder}-- A chronic mood disorder characterized by numerous periods of hypomanic symptoms and periods of depressive symptoms over at least two years (one year in children and adolescents). The hypomanic and depressive symptoms do not meet criteria for a hypomanic or major depressive episode. The patient is never symptom-free for more than two months at a time. It carries a 15--50\% risk of developing bipolar I or II disorder.

\medterm{Defense Mechanisms}-- Unconscious psychological strategies used to manage anxiety and protect the ego from distressing thoughts or feelings. Mature defenses include sublimation, humor, and altruism. Immature defenses include projection, denial, splitting, and acting out. Understanding defense mechanisms helps clinicians conceptualize personality organization.

\medterm{Delayed Ejaculation}-- A sexual dysfunction characterized by a persistent or recurrent delay in, infrequency of, or absence of ejaculation during partnered sexual activity, lasting \(\ge\)6 months and causing marked distress. Prevalence: 1-5\% of men. Lifelong vs. acquired. Causes: psychosocial (masturbatory style discrepancy, relationship issues, performance pressure, religious upbringing), medication-induced (SSRIs, antipsychotics, alpha-blockers, thiazide diuretics), neurological (spinal cord injury, multiple sclerosis, diabetic neuropathy), surgical (radical prostatectomy, retroperitoneal lymph node dissection), and hormonal (low testosterone, hyperprolactinemia). Diagnosis: clinical interview distinguishing anejaculation from retrograde ejaculation (post-orgasmic urine analysis). Treatment: address underlying cause, switch/reduce offending medications, vibratory stimulation (for neurogenic), electroejaculation (for fertility), psychotherapy, and off-label cabergoline or oxytocin (limited evidence).

\medterm{Delirium}-- An acute, fluctuating disturbance of attention and awareness with an underlying medical cause. It is characterized by disorganized thinking, altered level of consciousness, and perceptual disturbances. Risk factors include advanced age, dementia, infection, medications, and metabolic derangements. Prevention through hydration, sleep promotion, and minimizing restraints is essential.

\medterm{Delirium Tremens}-- The most severe form of alcohol withdrawal syndrome, a medical emergency with mortality of 5--15\% if untreated. Typically occurs 48--96 hours after the last drink in individuals with chronic heavy alcohol use. Characterized by: altered mental status (confusion, disorientation), profound autonomic hyperactivity (tachycardia, hypertension, hyperthermia, diaphoresis), tremor, hallucinations (visual, tactile, auditory), agitation, and seizures. Risk factors: prior DTs, concurrent infection or illness, electrolyte abnormalities. Management: ICU admission, benzodiazepines (diazepam or lorazepam, high-dose titrated to symptom control), thiamine, correction of electrolyte disturbances, and supportive care (hydration, nutrition, monitoring).

\medterm{Delusional Disorder}-- A disorder characterized by the presence of one or more delusions lasting one month or more in the context of relatively preserved functioning. The delusions are non-bizarre (involving real-life situations that could conceivably occur, such as being followed, infected, loved from afar, or conspired against). If hallucinations are present, they are not prominent and are related to the delusional theme. The behavior is not obviously bizarre or odd beyond the impact of the delusion.

\medterm{Denial}-- A defense mechanism in which a person refuses to accept an unpleasant reality. It is a common initial response to traumatic events and terminal illness diagnoses. While adaptive in the short term, chronic denial can prevent appropriate medical treatment and psychological processing.

\medterm{Dependent Personality Disorder}-- A pattern of submissive and clinging behavior related to an excessive need to be taken care of. At least five of eight criteria: difficulty making everyday decisions without excessive advice, needing others to assume responsibility, difficulty expressing disagreement due to fear of loss of support, difficulty initiating projects, going to excessive lengths to obtain nurturance, feeling uncomfortable or helpless when alone, urgently seeking another relationship when one ends, and unrealistically preoccupied with fears of being left to care for self. Diagnosis is clinical using DSM-5 criteria. Treatment includes psychodynamic psychotherapy and CBT focusing on building autonomy and assertiveness. Pharmacotherapy is used for comorbid depression or anxiety. Prognosis is variable, with some improvement possible with long-term therapy.

\medterm{Depersonalization-Derealization Disorder}-- Persistent or recurrent experiences of depersonalization, derealization, or both. Depersonalization includes experiences of unreality, detachment, or being an outside observer with respect to one's thoughts, feelings, sensations, body, or actions. Derealization includes experiences of unreality or detachment with respect to surroundings. The individual maintains intact reality testing, meaning they recognize the experiences as subjective and not objectively real.

\medterm{Depression}-- Major depressive disorder is a common, disabling mood disorder characterized by persistent sadness and loss of interest or pleasure in activities, affecting approximately 7 percent of adults worldwide annually. According to DSM-5 criteria, diagnosis requires five or more of the following symptoms present during the same 2-week period representing a change from previous functioning, with at least one being either depressed mood or anhedonia: depressed mood most of the day nearly every day, marked diminished interest or pleasure in all or almost all activities most of the day nearly every day, significant weight loss or gain, insomnia or hypersomnia, psychomotor agitation or retardation, fatigue or loss of energy, feelings of worthlessness or excessive or inappropriate guilt, diminished ability to think or concentrate, and recurrent thoughts of death or suicidal ideation. The episode causes clinically significant distress or impairment. Depression is the leading cause of disability worldwide and is a major risk factor for suicide. Treatment includes antidepressants (SSRIs, SNRIs, bupropion, mirtazapine as first-line), psychotherapy (CBT, interpersonal therapy, behavioural activation), and combination therapy for moderate-to-severe depression. Electroconvulsive therapy (ECT) is the most effective treatment for severe depression with psychotic features or treatment resistance. Prognosis is variable; many patients achieve remission, but recurrence rates are high, particularly without maintenance treatment. Management of suicide risk is a priority.

\medterm{Dialectical Behavior Therapy}-- A comprehensive cognitive-behavioral treatment originally developed for borderline personality disorder. It combines cognitive-behavioral techniques with mindfulness and acceptance strategies. Core modules include mindfulness, distress tolerance, emotion regulation, and interpersonal effectiveness. DBT is the treatment of choice for borderline personality disorder and has been adapted for substance use disorders, eating disorders, and PTSD.

\medterm{Dissociation}-- A disruption in the normal integration of consciousness, memory, identity, emotion, perception, behavior, and sense of self. It exists on a spectrum from normal (daydreaming) to pathological (DID, depersonalization-derealization). It is strongly associated with childhood trauma and can be triggered by stress, sleep deprivation, and certain substances.

\medterm{Dissociative Amnesia}-- Inability to recall important autobiographical information, usually of a traumatic or stressful nature, that is inconsistent with ordinary forgetting. The amnesia may be localized, selective, generalized, or continuous. In dissociative amnesia with dissociative fugue, there is purposeful travel or bewildered wandering associated with amnesia for identity or other important autobiographical information.

\medterm{Dissociative Identity Disorder}-- The presence of two or more distinct personality states or an experience of possession involving marked discontinuity in sense of self, agency, affect, behavior, memory, perception, cognition, and/or sensory-motor functioning. Identity disruption is accompanied by gaps in recall of everyday events, personal information, and/or traumatic events that are inconsistent with ordinary forgetting. The disturbance causes clinically significant distress or impairment and is not part of a broadly accepted cultural or religious practice. Caused by severe childhood trauma, particularly prolonged and severe abuse in early childhood. The presentation is often misunderstood and can be confused with malingering. Diagnosis requires careful clinical assessment using structured interviews (SCID-D). Treatment focuses on trauma-informed psychotherapy, with phases of stabilisation, trauma processing, and integration. Medications treat comorbid symptoms. Prognosis is guarded.

\medterm{Dysthymia (Persistent Depressive Disorder)}-- A chronic form of depression with depressed mood for most of the day, more days than not, for \(\ge\)2 years (\(\ge\)1 year in children/adolescents). Symptoms: poor appetite or overeating, insomnia or hypersomnia, low energy, low self-esteem, poor concentration, and hopelessness. Does not meet full criteria for major depressive episode continuously. Double depression: MDE superimposed on dysthymia. Treatment: SSRIs/SNRIs and psychotherapy (CBT, interpersonal therapy).

\medterm{Eating Disorder}-- A group of conditions characterized by abnormal eating habits and distorted body image, including anorexia nervosa, bulimia nervosa, binge eating disorder, and avoidant/restrictive food intake disorder (ARFID). Anorexia: severe calorie restriction, intense fear of weight gain, and body image disturbance with dangerously low weight. Bulimia: binge eating followed by purging behaviors. Binge eating: recurrent binge eating without purging. Treatment: nutritional rehabilitation, CBT, family-based therapy for adolescents, and fluoxetine for bulimia. Medical complications require monitoring.

\medterm{Electroconvulsive Therapy}-- A neuromodulation treatment involving the induction of a generalized seizure under general anesthesia using electrical current applied to the scalp. It is the most effective treatment for severe depression, especially with psychotic features, catatonia, or suicidal ideation. It is also effective for severe mania and treatment-resistant schizophrenia. Cognitive side effects include transient memory impairment. Bilateral electrode placement is more effective but has more cognitive side effects than unilateral placement. ECT is highly effective, with response rates of 70-90 percent for severe depression. The mechanism of action is not fully understood but involves modulation of neurotransmitter systems, neuroplasticity, and increased hippocampal neurogenesis. Modern ECT is performed under general anesthesia with muscle relaxation, using brief electrical pulses to induce a generalised seizure lasting approximately 20-60 seconds. Contraindications include raised intracranial pressure, recent myocardial infarction, and unstable vascular aneurysms. Memory loss (both retrograde and anterograde) is the most significant side effect and is usually transient. Maintenance ECT is sometimes used to prevent relapse.

\medterm{Erectile Disorder}-- A sexual dysfunction characterized by persistent or recurrent difficulty achieving or maintaining an erection sufficient for sexual activity, or a marked decrease in erectile rigidity, lasting \(\ge\)6 months and causing marked distress. Prevalence: 10-20\% of men under 60, rising to 40-70\% over age 70. Etiology: organic (vascular disease, diabetes, hypertension, hyperlipidemia, obesity, smoking, hypogonadism, neurological conditions, postsurgical) and psychogenic (performance anxiety, depression, relationship conflict). Diagnosis: clinical history, IIEF-5 questionnaire, nocturnal penile tumescence testing, penile duplex Doppler ultrasound, and hormonal profile. First-line treatment: PDE5 inhibitors (sildenafil, tadalafil, vardenafil, avanafil). Second-line: intracavernosal injections (alprostadil, papaverine, phentolamine), vacuum erection devices, and low-intensity shockwave therapy. Third-line: penile prosthesis surgery. Treat underlying medical conditions and offer psychotherapy for psychogenic factors.

\medterm{Euphoria}-- An intense feeling of well-being, elation, or happiness, which can be a normal emotional response or a pathologic symptom. In medical contexts, euphoria is most commonly associated with: substance use (opioids, stimulants--cocaine, amphetamines, MDMA, cannabis, alcohol), psychiatric disorders (mania in bipolar I disorder--elevated, expansive, or irritable mood with increased energy), neurologic conditions (multiple sclerosis--lesional euphoria, frontal lobe syndromes), and endocrine disorders (Cushing syndrome, hyperthyroidism). Euphoric mood is distinguished from euthymia (normal mood) and from the emotional lability of pseudobulbar affect. Pathologic euphoria may impair judgment and lead to risky behavior.

\medterm{Excoriation Disorder}-- Recurrent skin picking resulting in skin lesions. The individual has made repeated attempts to decrease or stop the skin picking. The skin picking causes clinically significant distress or impairment. It is not attributable to a medical condition or substance. Common sites include face, arms, and hands. The picking may be automatic or focused and is often preceded by tension and followed by relief.

\medterm{Exhibitionistic Disorder}-- A paraphilic disorder involving recurrent, intense sexual arousal from exposing one's genitals to an unsuspecting person. The individual has acted on these urges or they cause marked distress or interpersonal difficulty. Duration \(\ge\)6 months. Onset typically before age 18. Treatment: CBT and SSRIs.

\medterm{Extrapyramidal Symptoms}-- Drug-induced movement disorders caused by D2 receptor blockade in the nigrostriatal pathway. They include acute dystonia (muscle spasm, typically within hours of starting), akathisia (inner restlessness), drug-induced parkinsonism (tremor, rigidity, bradykinesia after weeks of use), and tardive dyskinesia (involuntary movements after months to years). EPS is more common with high-potency typical antipsychotics. Anticholinergic agents (benztropine, diphenhydramine) treat acute dystonia and parkinsonism, and beta-blockers (propranolol) treat akathisia. Tardive dyskinesia is managed by switching to a second-generation antipsychotic with lower D2 potency (clozapine, quetiapine) and adding VMAT2 inhibitors (valbenazine, deutetrabenazine). Prevention involves using the lowest effective dose of antipsychotic and regular monitoring using the Abnormal Involuntary Movement Scale (AIMS).

\medterm{Factitious Disorder}-- Presentation of physical or psychological symptoms or imposing injury or disease on another, with the deceptive intention of assuming the sick role. In factitious disorder imposed on self, the individual falsifies signs or symptoms, induces injury or disease, and presents to medical providers. In factitious disorder imposed on another, the individual falsifies symptoms in another person. The deceptive behavior occurs even in the absence of obvious external rewards.

\medterm{Female Orgasmic Disorder (Anorgasmia)}-- A sexual dysfunction characterized by a persistent or recurrent delay in, infrequency of, or absence of orgasm, or markedly reduced intensity of orgasmic sensation, in \(\ge\)75\% of sexual encounters, lasting \(\ge\)6 months and causing marked distress. Lifelong (primary, never achieved orgasm by any means) vs. acquired (previously able) vs. situational (only with certain partners or activities). Prevalence: 10-15\% of women. Causes: psychosocial (anxiety, distraction, lack of knowledge, relationship issues, negative attitudes toward sex), medication-induced (SSRIs — most common, antipsychotics), medical (spinal cord injury, multiple sclerosis, diabetes, pelvic surgery), and partner factors. Diagnosis: clinical DSM-5 criteria. Treatment: education, directed masturbation, CBT, mindfulness-based sex therapy, sensate focus, and vibratory stimulation. Switching to bupropion or adding buspirone may help with SSRI-induced anorgasmia. Filbanserin and bremelanotide are approved for HSDD not anorgasmia specifically.

\medterm{Female Sexual Interest/Arousal Disorder}-- A sexual dysfunction in women characterized by absent or reduced sexual interest, thoughts, or receptivity, and/or absent or reduced sexual arousal (genital or nongenital), persisting for \(\ge\)6 months and causing marked distress. Combines the former DSM-IV diagnoses of hypoactive sexual desire disorder and female sexual arousal disorder. Prevalence: 10-20\% of women. Etiology: multifactorial (relationship factors, partner factors, individual vulnerability, medical/psychiatric comorbidities, medication effects). Diagnosis: clinical DSM-5 criteria. Treatment: psychotherapy (CBT, mindfulness-based sex therapy), sensate focus exercises, education, addressing relationship factors, and pharmacotherapy (flibanserin, bremelanotide, off-label testosterone therapy). Lubricants and vaginal moisturizers for arousal complaints.

\medterm{Fetishistic Disorder}-- A paraphilic disorder characterized by recurrent, intense sexual arousal from either the use of nonliving objects (fetish) or a highly specific focus on nongenital body parts (partialism), lasting \(\ge\)6 months and causing distress or impairment or involving acting on these urges. Common fetishes: shoes, feet, footwear, leather, lingerie, and rubber. Fetishism is common in the general population (up to 30\% report some fetish interest) and typically only becomes a disorder when it causes significant distress, impairment, or involves nonconsenting individuals or causes harm. Most common in males. Treatment: CBT, aversion therapy (historical, controversial), and SSRIs.

\medterm{Forensic Psychiatry}-- The interface of psychiatry and the legal system. It involves evaluations for criminal responsibility, competency to stand trial, malingering assessment, risk assessment for violence, and civil commitment. Forensic psychiatrists serve as expert witnesses and provide opinions on legal questions. Malingering is the intentional production of false or grossly exaggerated symptoms for external gain.

\medterm{Frotteuristic Disorder}-- A paraphilic disorder involving recurrent, intense sexual arousal from touching or rubbing against a nonconsenting person, lasting \(\ge\)6 months and causing distress or impairment or involving acting on these urges. Typically occurs in crowded public places (subways, buses, concerts). The individual rubs their genitals against the victim's buttocks or thighs. Onset: adolescence, almost exclusively male. Typically decreasing with age. Legal consequences: sexual offense charges. Treatment: CBT, relapse prevention, and SSRIs.

\medterm{Gender Dysphoria}-- Psychological distress resulting from incongruence between one's experienced gender and assigned gender at birth. Persists for \(\ge\)6 months and causes significant impairment. Onset in childhood (may resolve) or adolescence/adulthood (persistent). Diagnosis: DSM-5 requires marked incongruence between experienced/expressed gender and primary/secondary sex characteristics, with strong desire to be rid of them or to have opposite sex characteristics. Treatment: gender-affirming care including hormone therapy (testosterone/estrogen), puberty blockers for adolescents, and surgical interventions. Access to affirming care improves mental health outcomes.

\medterm{Generalized Anxiety Disorder}-- Characterized by excessive, difficult-to-control anxiety and worry about multiple events or activities occurring more days than not for at least six months. Associated symptoms include restlessness, being easily fatigued, difficulty concentrating, irritability, muscle tension, and sleep disturbance. The anxiety causes significant distress or impairment. It is often comorbid with depression and other anxiety disorders.

\medterm{Genito-Pelvic Pain/Penetration Disorder}-- A sexual dysfunction characterized by persistent or recurrent difficulties with vaginal penetration, marked vulvovaginal or pelvic pain during intercourse, fear or anxiety about penetration, and/or tensing of the pelvic floor muscles, lasting \(\ge\)6 months and causing marked distress. Combines the former DSM-IV diagnoses of vaginismus and dyspareunia. Prevalence: 3-15\% of women. Etiology: multifactorial (pelvic floor hypertonicity, history of sexual trauma, anxiety, endometriosis, vulvodynia, interstitial cystitis, postpartum/perimenopausal changes, inadequate lubrication). Diagnosis: clinical interview, pelvic exam, and exclusion of organic causes (infection, endometriosis, adhesions). Treatment: multimodal — pelvic floor physical therapy, dilator therapy, CBT/sex therapy, lubricants, topical lidocaine, vaginal estrogen (if atrophic), trigger point injections, and treatment of underlying medical conditions. Graded exposure to penetration is key.

\medterm{Geriatric Psychiatry}-- The subspecialty of psychiatry focused on the prevention, evaluation, diagnosis, and treatment of mental and emotional disorders in older adults. Geriatric patients often present atypically, with depression manifesting as somatic complaints and cognitive decline. Common conditions include late-life depression, dementia, delirium, and late-onset psychosis. Polypharmacy, medical comorbidity, and altered drug metabolism complicate treatment.

\medterm{Grief}-- A natural emotional response to loss, particularly the death of a loved one, but also following other significant losses such as divorce, job loss, or declining health. It encompasses a range of feelings including sadness, anger, guilt, and numbness, often occurring in waves rather than linear stages. While grief is not a mental disorder, prolonged or disabling grief may be diagnosed as prolonged grief disorder and benefit from specialized therapy.

\medterm{Hallucination}-- A perception in the absence of an external stimulus, experienced as real and occurring in objective space. Hallucinations can involve any sensory modality: auditory (most common, voices commenting or conversing--characteristic of schizophrenia, also in bipolar mania, psychotic depression), visual (formed or unformed--delirium, dementia, Charles Bonnet syndrome in visually impaired with preserved cognition, Lewy body dementia), olfactory (unusual odours--temporal lobe epilepsy, olfactory reference syndrome), gustatory (tastes--seizures, psychosis), tactile (formication--cocaine/amphetamine use, alcohol withdrawal, delusional infestation), and somatic (sensations within the body). Hallucinations are distinct from: illusions (distortions of real external stimuli), hallucinations in sleep (hypnagogic at sleep onset, hypnopompic on waking--normal if not distressing), and pseudohallucinations (experienced as lacking objective reality, located inside the head, often in PTSD). Hallucinations in the context of psychiatric illness must be differentiated from those due to: delirium (acute, fluctuating, medical cause), dementia (chronic, progressive), epilepsy (ictal/postictal, especially temporal lobe), drugs (stimulants, hallucinogens like LSD, psilocybin, ketamine, withdrawal from alcohol/benzodiazepines), sensory deprivation, and migraine (visual aura). Management: treat underlying cause; antipsychotics for psychiatric disorders (olanzapine, risperidone, clozapine for refractory auditory hallucinations); ECT for severe psychotic depression.

\medterm{Histrionic Personality Disorder}-- A pattern of excessive emotionality and attention-seeking. At least five of eight criteria: discomfort when not the center of attention, inappropriate sexually seductive behavior, displays rapidly shifting and shallow emotions, uses physical appearance for attention, speaks in an impressionistic and vague manner, shows self-dramatization and theatricality, is suggestible, and considers relationships more intimate than they actually are.

\medterm{Hoarding Disorder}-- Persistent difficulty discarding or parting with possessions, regardless of their actual value. The individual perceives a need to save items and experiences distress at the thought of discarding them. This results in accumulation of possessions that congest and clutter active living areas. The hoarding is not attributable to another medical condition and causes significant distress or impairment. It is distinct from collecting, which is organized and purposeful.

\medterm{Hypoactive Sexual Desire Disorder (HSDD)}-- A sexual dysfunction characterized by persistently or recurrently deficient or absent sexual thoughts, fantasies, and desire for sexual activity, causing marked distress or interpersonal difficulty. In DSM-5, HSDD in women is merged with female sexual arousal disorder into female sexual interest/arousal disorder. In men, it remains a separate diagnosis. Prevalence: 10-15\% of men, 20-30\% of women. Causes: psychosocial (relationship conflict, stress, depression, anxiety), hormonal (low testosterone, hyperprolactinemia, menopause), medication-induced (SSRIs, antipsychotics, hormonal contraceptives), and medical (diabetes, thyroid disease, chronic illness). Diagnosis: clinical interview ruling out other psychiatric disorders, medication effects, and medical conditions. Lab work: testosterone (total and free), prolactin, TSH. Treatment: treat underlying cause, modify offending medications, psychotherapy (CBT, sensate focus), and for women: flibanserin (Addyi) or bremelanotide (Vyleesi); for men with low testosterone: testosterone replacement therapy.

\medterm{Hypomanic Episode}-- A distinct period of abnormally elevated, expansive, or irritable mood with increased energy lasting at least four consecutive days. The episode is not severe enough to cause marked impairment in functioning or necessitate hospitalization, and there are no psychotic features. The change in functioning is observable by others. Hypomanic episodes are the defining feature of bipolar II disorder and may transition to full mania or major depression.

\medterm{Hysteria}-- A historical psychiatric term that has been discarded from modern diagnostic classifications (DSM-5, ICD-11). Originally derived from the Greek word for uterus (hystera), reflecting the ancient belief that symptoms originated from a 'wandering uterus'. Previously used to describe a range of symptoms: (1) Motor: paralysis, seizures, pseudoseizures (psychogenic non-epileptic seizures), mutism, aphonia; (2) Sensory: anesthesia, paraesthesia, blindness, deafness; (3) Dissociative: amnesia, fugue, depersonalisation, multiple personality; (4) Somatic: chronic pain, globus hystericus, hyperventilation. The term hysteria was replaced by more specific diagnostic categories: conversion disorder (functional neurological symptom disorder--motor/sensory symptoms without organic cause), somatisation disorder (now somatic symptom disorder, multiple physical complaints), dissociative disorders (dissociative identity disorder, dissociative amnesia). Freud redefined hysteria as resulting from repressed psychological trauma (especially sexual abuse) with unconscious conversion into physical symptoms. Pierre Janet proposed dissociation as the core mechanism. Current understanding integrates psychodynamic, behavioural, cognitive, and neurobiological factors, and does not attribute to the uterus.

\medterm{Ideas of Reference}-- The false belief that unrelated events, objects, or people have a particular and unusual significance specifically for the patient. For example, believing that TV newscasters are sending personal messages or that strangers are conspiring against them. While common in schizotypal personality disorder, they can also occur in schizophrenia and bipolar disorder.

\medterm{Illness Anxiety Disorder}-- Preoccupation with having or acquiring a serious illness, with minimal or no somatic symptoms present. The individual has a high level of anxiety about health, performs excessive health-related behaviors, and experiences maladaptive avoidance. Symptoms persist for at least six months. The individual is not delusional and insight may vary. Previously known as hypochondriasis.

\medterm{Insomnia Disorder}-- Difficulty initiating or maintaining sleep, or early-morning awakening with inability to return to sleep, at least three nights per week for at least three months. The individual experiences clinically significant distress or impairment in functioning. Despite adequate opportunity for sleep, the individual is unable to obtain sufficient sleep. It is not better explained by another sleep disorder, substance use, or medical condition. Cognitive behavioral therapy for insomnia (CBT-I) is first-line treatment. Pharmacological options include benzodiazepine receptor agonists (zolpidem, zaleplon, eszopiclone), melatonin receptor agonists (ramelteon), orexin receptor antagonists (suvorexant, lemborexant), and sedating antidepressants (trazodone, doxepin, mirtazapine). Benzodiazepines are generally avoided due to tolerance and dependence risks. Sleep hygiene measures and stimulus control are essential components of management.

\medterm{Intermittent Explosive Disorder}-- A disorder characterized by recurrent, impulsive, problematic outbursts of verbal or physical aggression that are grossly disproportionate to the provocation. The outbursts are not premeditated and cause marked distress or functional impairment. Age >6 years. Not better explained by another disorder (ADHD, conduct disorder, PTSD, ASD) or substance intoxication. Diagnosis: DSM-5 criteria including aggressive outbursts occurring twice weekly for 3 months (verbal/non-destructive) or 3 outbursts involving damage/injury in the past year. Treatment: CBT (anger management, relaxation), SSRIs (fluoxetine), and mood stabilizers.

\medterm{Interpersonal Psychotherapy}-- A time-limited, evidence-based psychotherapy for depression that focuses on the interpersonal context of mood symptoms. It identifies four problem areas: grief, role disputes, role transitions, and interpersonal deficits. IPT helps patients improve interpersonal functioning and social support. It is an alternative to CBT for depression and has been adapted for eating disorders, postpartum depression, and adolescent depression.

\medterm{Leukotomy}-- A neurosurgical procedure involving the cutting of white matter fibers in the prefrontal cortex, also known as frontal lobotomy or prefrontal leukotomy. Developed by Egas Moniz (Nobel Prize 1949) and Walter Freeman. leukotomy was used for severe psychiatric disorders (schizophrenia, severe depression, OCD) from the 1930s to 1950s, before the advent of effective psychopharmacology. The transorbital lobotomy (Freeman, using an ice pick-like instrument through the orbit) was a simplified outpatient procedure. Use declined dramatically in the 1950s with the introduction of antipsychotics (chlorpromazine) and antidepressants, and due to severe lasting side effects. Side effects of leukotomy: apathy, emotional blunting, inattention, executive dysfunction, cognitive impairment, seizures (5-10\%), and death. leukotomy is rarely performed today, but modern stereotactic procedures (anterior capsulotomy, cingulotomy, subcaudate tractotomy) are used in extremely severe, treatment-refractory cases of OCD and depression.

\medterm{Lithium}-- A mood-stabilising medication used primarily for bipolar disorder (mania, hypomania, and prophylaxis of recurrent episodes). See Lithium Toxicity and Bipolar Disorder.

\medterm{Lithium Toxicity}-- Occurs when lithium levels exceed the therapeutic range. Mild toxicity (1.5--2.0 mEq/L): nausea, vomiting, diarrhea, coarse tremor, muscle weakness. Moderate toxicity (2.0--2.5 mEq/L): confusion, ataxia, blurred vision, dysarthria. Severe toxicity (above 2.5 mEq/L): seizures, coma, cardiac arrhythmias, renal failure, and potentially death. Treatment includes stopping lithium, IV fluids, and hemodialysis for severe cases. NSAIDs, ACE inhibitors, and thiazide diuretics increase lithium levels.

\medterm{Lysergic Acid Diethylamide (LSD)}-- A potent hallucinogenic drug (psychedelic) classified as a Schedule 1 controlled substance (high abuse potential, no accepted medical use). LSD acts as a partial agonist at serotonin 5-HT2A receptors in the prefrontal cortex, leading to altered sensory perception, mood, and cognition. Effects: visual hallucinations (geometric patterns, intensified colours, trails), synaesthesia (cross-modal sensory experience, e.g., hearing colours), profound introspection, euphoria/dysphoria, altered sense of time, depersonalization, and mystical experiences. Onset: 30-60 minutes after oral ingestion. Peak: 2-4 hours. Duration: 8-12 hours. Dose: typical 50-200 microgram. Adverse effects: 'bad trip' (anxiety, panic, paranoia, fear of losing control), flashbacks (recurrence of drug effect weeks to months later, Hallucinogen Persisting Perception Disorder--HPPD), exacerbation of pre-existing psychiatric illness (psychosis, bipolar disorder), and acute toxicity (hyperthermia, hypertension, tachycardia, seizures, serotonin syndrome). LSD is not considered physically addictive (does not produce withdrawal or compulsive drug-seeking behaviour), but tolerance develops rapidly (cross-tolerance with other 5-HT2A agonists like psilocybin and mescaline). Medical research: renewed interest in using LSD and other psychedelics (psilocybin, MDMA, ayahuasca) for treatment-resistant depression, anxiety in terminal illness, PTSD, and substance use disorders in controlled clinical settings.

\medterm{Major Depressive Disorder}-- A mood disorder characterized by at least one major depressive episode consisting of five or more symptoms during the same two-week period. At least one symptom must be either depressed mood or loss of interest or pleasure (anhedonia). Symptoms include significant weight change, insomnia or hypersomnia, psychomotor agitation or retardation, fatigue, feelings of worthlessness or excessive guilt, diminished concentration, and recurrent thoughts of death or suicidal ideation. The episode causes clinically significant distress or impairment. MDD is the leading cause of disability worldwide, affecting approximately 7 percent of adults annually. Treatment involves antidepressants (SSRIs, SNRIs, bupropion, mirtazapine), psychotherapy (CBT, interpersonal therapy, behavioural activation), and combination therapy for moderate-to-severe episodes. ECT is reserved for severe, psychotic, or treatment-resistant depression. Prognosis is variable; recurrence is common, and maintenance treatment is recommended.

\medterm{Male Hypoactive Sexual Desire Disorder}-- A sexual dysfunction in men characterized by persistently deficient or absent sexual/erotic thoughts or desire for sexual activity, causing marked distress. Prevalence: 5-15\% of men. Distinguish from erectile disorder (men may have ED but intact desire). Causes: low testosterone (primary or secondary hypogonadism), hyperprolactinemia, depression, relationship conflict, medication side effects (SSRIs, antipsychotics, beta-blockers, spironolactone), chronic illness (diabetes, CKD, COPD), and aging. Diagnosis: total and free testosterone, prolactin, TSH, and clinical interview. Treatment: testosterone replacement (if confirmed hypogonadal), treat hyperprolactinemia, switch offending medications, CBT, and sex therapy. Testosterone is not recommended for age-related low T without clear hypogonadism.

\medterm{Malingering}-- The intentional fabrication or exaggeration of physical or psychological symptoms for external incentives (avoiding work, obtaining financial compensation, evading criminal prosecution, obtaining drugs). Not a psychiatric illness (coded as a V-code/Z-code). Suspicion raised by: non-compliance with treatment, antisocial personality disorder, discrepancies between reported symptoms and objective findings, and symptoms that emerge in specific contexts (legal, military, disability evaluation). Diagnosis: collateral information, psychological testing (MMPI-2 validity scales). No psychiatric treatment; management involves confrontation and documentation.

\medterm{Mania}-- A distinct period of abnormally elevated, expansive, or irritable mood with increased goal-directed activity or energy, lasting \(\ge\)1 week (or any duration if hospitalization required). Requires \(\ge\)3 of: inflated self-esteem, decreased need for sleep, pressured speech, flight of ideas, distractibility, increased goal-directed activity, and excessive involvement in risky activities. Causes significant functional impairment or psychosis. Core feature of bipolar I disorder. Treatment: mood stabilizers (lithium, valproate) and atypical antipsychotics.

\medterm{Manic Episode}-- A distinct period lasting at least seven days (or any duration if hospitalization is required) of abnormally elevated, expansive, or irritable mood with increased energy or activity. At least three symptoms must be present: inflated self-esteem or grandiosity, decreased need for sleep, pressured speech, flight of ideas or racing thoughts, distractibility, increased goal-directed activity or psychomotor agitation, and excessive involvement in pleasurable activities with high potential for painful consequences (e.g., buying sprees, sexual indiscretions, risky investments). The mood disturbance is severe enough to cause marked impairment in social or occupational functioning, requires hospitalisation to prevent harm to self or others, or includes psychotic features. First-line treatment is lithium or valproate, with atypical antipsychotics (olanzapine, quetiapine, risperidone, aripiprazole, and asenapine) as effective alternatives. Benzodiazepines are used as adjunctive therapy for acute agitation. ECT is highly effective for severe or refractory mania.

\medterm{Melancholic Features}-- A specifier for major depressive episodes characterized by a loss of interest or pleasure in all or almost all activities, a depressed mood that is qualitatively distinct from grief, prominent psychomotor retardation or agitation, significant anorexia or weight loss, excessive guilt, and worsening of mood in the morning. It suggests a more severe, biologically driven depression with better response to antidepressants or ECT.

\medterm{Melatonin Receptor Agonists}-- Medications that act on MT1 and MT2 receptors in the suprachiasmatic nucleus to regulate circadian rhythms. Ramelteon is approved for insomnia with sleep onset difficulty. It has no abuse potential and does not cause dependence. It is a useful alternative to benzodiazepines for sleep initiation problems.

\medterm{Mental Health}-- Mental health encompasses emotional, psychologic, and social well-being, affecting how individuals think, feel, behave, and handle stress, relationships, and decision-making. Common mental disorders: anxiety disorders (generalized anxiety, panic, phobias, social anxiety—lifetime prevalence ~29\%), depressive disorders (major depressive disorder, persistent depressive disorder—lifetime prevalence ~21\%), bipolar disorder (manic-depressive illness), schizophrenia and other psychotic disorders, and trauma/stressor-related disorders (PTSD, adjustment disorder). The biopsychosocial model recognizes genetic vulnerability (heritability 30--50\% for major depression), neurobiologic factors (neurotransmitter dysregulation, HPA-axis abnormalities), psychological factors (maladaptive cognitions, coping deficits), and social determinants (childhood adversity, poverty, discrimination, isolation). Treatment: psychotherapy (CBT, DBT, psychodynamic, interpersonal therapy), pharmacotherapy (antidepressants, mood stabilizers, antipsychotics, anxiolytics), neuromodulation (ECT, TMS, ketamine/esketamine for treatment-resistant depression), and social support interventions. Screening (PHQ-9, GAD-7, AUDIT-C) in primary care improves detection and outcomes.

\medterm{Mental Status Examination}-- A structured assessment of the patient's current psychological functioning performed during the clinical interview. It evaluates appearance, behaviour, speech, mood, affect, thought process, thought content, perceptions, cognition, insight, and judgment. It provides a snapshot of the patient's mental state at the time of evaluation and is documented in the psychiatric note.

\medterm{Methylphenidate}-- See Methylphenidate in Pharmacology.

\medterm{Monoamine Oxidase Inhibitors}-- Antidepressants that inhibit monoamine oxidase, the enzyme responsible for breaking down serotonin, norepinephrine, and dopamine. MAOIs include phenelzine, tranylcypromine, isocarboxazid, and selegiline (transdermal). They are effective for atypical depression and treatment-resistant depression. The major risk is hypertensive crisis from tyramine-rich foods (``cheese reaction''). They are contraindicated with sympathomimetics, meperidine, and other serotonergic agents due to risk of serotonin syndrome. Dietary restrictions require avoidance of aged cheeses, cured meats, fermented foods, and certain alcoholic beverages (tyramine-rich foods). The transdermal selegiline patch bypasses gastrointestinal MAO and does not require dietary restrictions at the lowest dose. MAOIs are reserved for treatment-resistant cases due to their side effect profile and drug interactions.

\medterm{Mood Stabilizers}-- Medications used to treat mood episodes in bipolar disorder and prevent recurrence. The primary mood stabilizers are lithium, valproate, carbamazepine, and lamotrigine. Each has a different mechanism and side effect profile. Lithium is the only agent with demonstrated anti-suicide properties. Valproate is particularly effective for rapid-cycling bipolar disorder. Lamotrigine is more effective for bipolar depression than mania.

\medterm{Motivational Interviewing}-- A client-centered, directive counseling method for enhancing intrinsic motivation to change by exploring and resolving ambivalence. Developed by Miller and Rollnick, it uses the spirit of partnership, acceptance, compassion, and evocation. Core skills include open-ended questions, affirmations, reflections, and summaries. It is particularly effective for substance use disorders and health behavior change.

\medterm{Narcissistic Personality Disorder}-- A pattern of grandiosity, need for admiration, and lack of empathy. At least five of nine criteria: grandiose sense of self-importance, preoccupation with fantasies of unlimited success, belief of being ``special,'' requires excessive admiration, sense of entitlement, interpersonally exploitative, lacks empathy, envious of others or believes others are envious, and shows arrogant, haughty behaviors. May present as grandiose or vulnerable (covert) narcissism.

\medterm{Narcolepsy}-- A chronic sleep disorder characterized by excessive daytime sleepiness and irresistible sleep attacks, with or without cataplexy (sudden bilateral loss of muscle tone triggered by strong emotions). Narcolepsy type 1: with cataplexy, caused by loss of hypocretin/orexin neurons in the hypothalamus. Other features: sleep paralysis (inability to move upon falling asleep or waking), hypnagogic/hypnopompic hallucinations, and disrupted nighttime sleep. Diagnosis: polysomnography followed by MSLT (mean sleep latency <8 minutes, \(\ge\)2 SOREM periods). Treatment: wake-promoting agents (modafinil, armodafinil, solriamfetol), sodium oxybate for cataplexy/EDS, and SSRIs/SNRIs for cataplexy.

\medterm{Neuroleptic}-- An older term for antipsychotic medications, derived from their ability to produce a state of mental indifference and motor slowing. The term is now largely replaced by ``antipsychotic,'' though it is still used in specific contexts such as neuroleptic malignant syndrome. Typical neuroleptics are first-generation antipsychotics.

\medterm{Neuroplasticity}-- The brain's ability to reorganize its structure, function, and connections in response to experience, learning, and injury. It is relevant to psychiatric treatment because antidepressants, psychotherapy, and ECT are thought to promote neuroplastic changes including neurogenesis in the hippocampus and synaptic remodeling.

\medterm{Obsessions}-- Recurrent, persistent, intrusive, and unwanted thoughts, urges, or images that are experienced as intrusive and cause marked anxiety or distress. Common obsessions include contamination fears, fears of harming oneself or others, need for symmetry, taboo thoughts (sexual, religious, or aggressive), and doubt. The individual recognizes the obsessions as excessive or irrational in many cases (good insight).

\medterm{Obsessive-Compulsive Disorder}-- Characterized by the presence of obsessions, compulsions, or both. Obsessions are recurrent, persistent, intrusive, and unwanted thoughts, urges, or images that cause marked anxiety or distress. The individual attempts to ignore or suppress them or neutralize them with some other thought or action. Compulsions are repetitive behaviors or mental acts that the individual feels driven to perform in response to an obsession or according to rigid rules. They are aimed at preventing or reducing anxiety or preventing a dreaded event, but are not connected in a realistic way with what they are designed to neutralise. Compulsions are time-consuming (taking more than one hour per day) and cause significant distress or impairment. Insight varies: good or fair insight (the individual recognises the beliefs are probably not true), poor insight, or absent insight (delusional beliefs). First-line treatment is CBT with exposure and response prevention (ERP) and SSRIs (fluoxetine, fluvoxamine, sertraline, paroxetine, or escitalopram) at higher doses than used for depression. The combination of CBT and SSRI is superior to either alone. Prognosis is variable; symptoms may be chronic but are usually responsive to treatment.

\medterm{Obsessive-Compulsive Personality Disorder}-- A pattern of preoccupation with orderliness, perfectionism, and mental and interpersonal control, at the expense of flexibility and efficiency. At least four of eight criteria: preoccupied with details, rules, lists, organization, or schedules to the extent that the major point of the activity is lost; shows perfectionism that interferes with task completion; excessively devoted to work; overconscientious and inflexible about morality; unable to discard worn-out objects; reluctant to delegate tasks unless others submit to their way of doing things. Diagnosis is clinical using DSM-5 criteria. Treatment is challenging as patients rarely seek help voluntarily. Psychodynamic therapy and CBT may help address perfectionism and control. Pharmacotherapy targets comorbid depression or anxiety. Prognosis is guarded, as the personality traits are ego-syntonic and resistant to change.

\medterm{OCD}-- See Obsessive-Compulsive Disorder.

\medterm{Opioid Overdose}-- A medical emergency characterized by the classic triad of respiratory depression, miosis (pinpoint pupils), and decreased level of consciousness. It can progress to respiratory arrest, cardiovascular collapse, and death. Treatment involves airway management, naloxone administration (0.4--2 mg IV/IM, may repeat), and monitoring for re-narcotization as naloxone wears off. Fentanyl and its analogs require higher and repeated naloxone doses.

\medterm{Opioid Use Disorder}-- A problematic pattern of opioid use leading to significant impairment or distress. Opioids include heroin, fentanyl, morphine, oxycodone, and hydrocodone. Criteria are the same eleven as for substance use disorders generally. Overdose presents with respiratory depression, miosis, and decreased level of consciousness. Naloxone is the antidote. Treatment includes methadone, buprenorphine, and naltrexone.

\medterm{Oppositional Defiant Disorder}-- A childhood/adolescent disorder characterized by a persistent pattern (\(\ge\)6 months) of angry/irritable mood, argumentative/defiant behavior, and vindictiveness. Symptoms: often loses temper, is touchy/easily annoyed, angry/resentful, argues with authority figures, actively defies rules, deliberately annoys others, blames others for mistakes, and has been spiteful/vindictive at least twice in the past 6 months. Milder than conduct disorder (no aggression toward people/animals, no destruction of property). Risk factor for development of conduct disorder. Treatment: parent management training, CBT, and school-based interventions.

\medterm{Orexin Receptor Antagonists}-- Medications that block orexin (hypocretin) receptors to promote sleep. Suvorexant and lemborexant are approved for insomnia. They reduce wakefulness by blocking the arousal system. They have lower abuse potential than benzodiazepines and do not cause respiratory depression.

\medterm{OSFED}-- Other Specified Feeding or Eating Disorder, used when the clinical presentation causes significant distress or impairment but does not meet full criteria for any of the specific eating disorders. Examples include atypical anorexia nervosa (all criteria met except weight is within or above normal range), subthreshold bulimia nervosa or binge eating disorder, purging disorder, and night eating syndrome.

\medterm{Other Mental Health Disorders}-- conditions beyond depression and anxiety, such as bipolar disorder (manic-depressive illness cycling between mania/hypomania and depression; lithium is cornerstone therapy), schizophrenia and psychotic disorders (hallucinations, delusions, disorganized thinking; dopamine hypothesis; antipsychotic treatment), obsessive-compulsive disorder (OCD; intrusive thoughts and repetitive rituals; CBT/ERP and SSRIs), post-traumatic stress disorder (PTSD; re-experiencing, avoidance, hyperarousal after trauma; trauma-focused therapy, SSRIs, prazosin for nightmares), eating disorders (anorexia nervosa, bulimia nervosa, binge-eating disorder; requiring medical monitoring for refeeding syndrome), personality disorders (borderline, antisocial, avoidant; dialectical behavior therapy for BPD), and substance use disorders (alcohol, opioids, stimulants; pharmacotherapy—naltrexone, buprenorphine—and psychosocial support). Multidisciplinary treatment: psychotherapy, pharmacotherapy, social support, and coordinated care with medical providers for substance-related and other medical complications.

\medterm{Other Specified Sexual Dysfunction}-- A category for sexual dysfunction symptoms that cause distress but do not meet full criteria for any specific sexual dysfunction. Examples: sexual aversion (extreme avoidance of or aversion to genital sexual contact), inadequate vaginal lubrication without pain, persistent genital arousal disorder (unwanted genital arousal without sexual desire, PGAD), and sexual dysfunction due to medical conditions not otherwise covered. Management is tailored to the specific presentation.

\medterm{Panic Attack}-- An abrupt surge of intense fear or discomfort reaching a peak within minutes. The individual experiences at least four of thirteen symptoms: palpitations, sweating, trembling, shortness of breath, feelings of choking, chest pain, nausea, dizziness, chills or heat, paresthesias, derealization or depersonalization, fear of losing control, and fear of dying. Panic attacks are not limited to panic disorder and occur in many psychiatric and medical conditions.

\medterm{Panic Disorder}-- Recurrent unexpected panic attacks, which are abrupt surges of intense fear or discomfort peaking within minutes. During attacks, individuals experience at least four of thirteen symptoms: palpitations, sweating, trembling, shortness of breath, feelings of choking, chest pain, nausea, dizziness, chills or heat sensations, paresthesias, derealization, fear of losing control, and fear of dying. At least one attack is followed by persistent worry about additional attacks or maladaptive behavioral changes aimed at avoiding future attacks. Panic disorder is a chronic condition with a relapsing-remitting course. The lifetime prevalence is approximately 3-5 percent, with onset typically in late adolescence or early adulthood. The fear of having another panic attack leads to a cycle of escalating anxiety and avoidance. Agoraphobia frequently co-occurs. Diagnosis is clinical using DSM-5 criteria. First-line treatment includes CBT with interoceptive exposure and SSRIs (sertraline, paroxetine, fluoxetine, escitalopram) or SNRIs (venlafaxine). Benzodiazepines (alprazolam, clonazepam) are effective acutely but avoided as first-line due to tolerance and dependence. Prognosis is good with treatment; most patients achieve significant symptom reduction.

\medterm{Paraphilia}-- A condition characterized by intense and persistent sexual interest in atypical objects, situations, or individuals. Examples: exhibitionism, fetishism, frotteurism, pedophilic disorder, sexual masochism, sexual sadism, transvestic disorder, voyeurism, and paraphilia not otherwise specified. Diagnosis requires that the paraphilia causes distress or impairment or involves harm to others. Treatment: CBT, SSRI medications, and antiandrogens in selected cases.

\medterm{Pedophilic Disorder}-- A paraphilic disorder characterized by intense and recurrent sexual attraction to prepubescent children (generally age \(\le\)13), lasting \(\ge\)6 months, and causing distress or impairment or involving acting on these urges. Diagnosis requires the individual is \(\ge\)16 years old and \(\ge\)5 years older than the child. Assessment: clinical interview, phallometric testing. Treatment: CBT focusing on relapse prevention and cognitive distortions; pharmacotherapy (SSRIs, antiandrogens — medroxyprogesterone acetate, cyproterone acetate, GnRH agonists like leuprolide) to reduce libido. Ethical and legal considerations are paramount.

\medterm{Perinatal Psychiatry}-- The subspecialty focused on psychiatric disorders during pregnancy and the postpartum period. Common conditions include perinatal depression, perinatal anxiety, perinatal OCD, and postpartum psychosis. Medication selection must balance maternal mental health with fetal or neonatal exposure risks. SSRIs are generally considered safe, while lithium and valproate carry teratogenic risks. Breastfeeding considerations influence medication choice.

\medterm{Persistent Depressive Disorder}-- formerly called dysthymia, a chronic depressive mood lasting at least two years in adults (one year in children and adolescents). Patients experience depressed mood for more days than not, along with at least two additional symptoms: poor appetite or overeating, insomnia or hypersomnia, low energy, low self-esteem, poor concentration, or feelings of hopelessness. During a two-year period, the patient has never been symptom-free for more than two months.

\medterm{Personality Disorder}-- A group of mental disorders characterized by enduring maladaptive patterns of thinking, feeling, and behaving that deviate from cultural expectations and cause functional impairment. Cluster A (odd/eccentric): paranoid, schizoid, schizotypal. Cluster B (dramatic/erratic): antisocial, borderline, histrionic, narcissistic. Cluster C (anxious/fearful): avoidant, dependent, obsessive-compulsive. Diagnosis requires onset in adolescence or early adulthood and stability over time. Treatment: psychotherapy (DBT for borderline, CBT, psychodynamic); limited pharmacotherapy for specific symptoms.

\medterm{Postpartum Depression}-- Major depressive or manic episodes occurring during pregnancy or within four weeks of delivery (DSM-5). The term is often used more broadly to cover depressive episodes up to one year postpartum. Risk factors include prior depression, lack of social support, stressful life events, and hormonal fluctuations. Postpartum psychosis is a distinct, more severe entity requiring immediate psychiatric hospitalization.

\medterm{Postpartum Psychosis}-- A psychiatric emergency occurring within the first two weeks postpartum, characterized by rapid onset of psychotic symptoms including confusion, disorientation, hallucinations, and delusions. It occurs in approximately 1--2 per 1000 deliveries and is associated with bipolar disorder. Risk of infanticide and suicide is elevated. Requires immediate psychiatric hospitalization and treatment with mood stabilizers and antipsychotics.

\medterm{Posttraumatic Stress Disorder (PTSD)}-- A psychiatric disorder that develops after exposure to a traumatic event (death, serious injury, sexual violence). Presents with 4 symptom clusters: re-experiencing (intrusive memories, flashbacks, nightmares), avoidance (of trauma reminders), negative alterations in cognition and mood, and hyperarousal (hypervigilance, exaggerated startle, irritability). Duration >1 month with functional impairment. Risk factors: prior trauma, lack of social support, female sex, and peritraumatic dissociation. Diagnosis: DSM-5 criteria. Treatment: trauma-focused CBT (prolonged exposure, cognitive processing therapy), EMDR, and SSRIs (sertraline, paroxetine). Prognosis: chronic without treatment.

\medterm{Premature (Early) Ejaculation}-- A sexual dysfunction characterized by a persistent or recurrent pattern of ejaculation occurring within approximately 1 minute of vaginal penetration and before the individual wishes it, occurring in \(\ge\)75\% of sexual encounters, lasting \(\ge\)6 months, and causing marked distress. Prevalence: 20-30\% of men, the most common male sexual dysfunction. Lifelong (primary) vs. acquired (secondary to ED, prostatitis, hyperthyroidism, anxiety, or relationship issues). Diagnosis: clinical DSM-5 criteria (IELT, perceived control, distress). Treatment: first-line — behavioral techniques (stop-start, squeeze technique), topical anesthetics (lidocaine-prilocaine cream). Pharmacotherapy: dapoxetine (short-acting SSRI, on-demand, approved in many countries outside the US), off-label SSRIs (paroxetine 10-40 mg, sertraline 50-200 mg) daily or on-demand, and clomipramine. Treat underlying ED if present.

\medterm{Premenstrual Dysphoric Disorder}-- A severe, debilitating form of premenstrual syndrome characterized by marked mood symptoms (depression, anxiety, irritability, mood swings), physical symptoms (breast tenderness, bloating, fatigue), and functional impairment that occur during the luteal phase and remit within days of menses onset. Symptoms must occur in most menstrual cycles for the past year and cause significant distress. Distinguishing from PMS: greater severity and predominance of affective symptoms. Diagnosis: prospective daily symptom charting for \(\ge\)2 cycles. Treatment: SSRIs (intermittent luteal phase or continuous), combined oral contraceptives (drospirenone-containing), CBT, and GnRH agonists in refractory cases.

\medterm{Projection}-- A defense mechanism in which unacceptable thoughts, feelings, or impulses are attributed to others. For example, a patient who is angry at the therapist may accuse the therapist of being angry at them. It is common in paranoid thinking, delusional disorders, and personality disorders.

\medterm{Psychodynamic Therapy}-- A form of depth psychology that explores how unconscious processes, early life experiences, and internal conflicts influence current thoughts, feelings, and behaviors. It involves exploration of transference and countertransference, defense mechanisms, and the therapeutic relationship as a vehicle for change. Brief psychodynamic therapy is time-limited and focused on specific conflicts. It is effective for depression, anxiety, personality disorders, and somatic symptom disorders.

\medterm{Psychogenic Non-Epileptic Seizures}-- Seizure-like episodes that are not caused by abnormal electrical activity in the brain. They are a form of conversion disorder with attacks or seizures. Events may include shaking, stiffening, unresponsiveness, and other movements. Diagnosis requires clinical evidence of incompatibility with epilepsy and can be confirmed by video-EEG monitoring showing normal EEG during an event. Treatment includes psychotherapy and education about the diagnosis.

\medterm{Psychopharmacology}-- The study and clinical application of medications used to treat psychiatric disorders. It involves understanding drug mechanisms, pharmacokinetics, drug interactions, and side effect profiles. It is the foundation of biological psychiatry and enables evidence-based medication management for mood, anxiety, psychotic, and cognitive disorders.

\medterm{PTSD}-- See Post-Traumatic Stress Disorder.

\medterm{Purging Disorder}-- Recurrent purging behavior to influence weight or shape in the absence of binge eating. Self-induced vomiting, misuse of laxatives, diuretics, or other medications. It is classified under Other Specified Feeding or Eating Disorder (OSFED). It can lead to electrolyte imbalances, dental erosion, Russell's sign, and esophageal tears.

\medterm{Reactive Attachment Disorder}-- A pattern of inhibited, emotionally withdrawn behavior toward adult caregivers. The child rarely or minimally seeks comfort when distressed and rarely or minimally responds to comfort when distressed. The child shows a pattern of persistently irritable and sad mood that is not limited to interactions with caregivers. It results from a pattern of insufficient care including social neglect, repeated changes of primary caregivers, or rearing in unusual settings.

\medterm{Restless Legs Syndrome}-- An urge to move the legs, usually accompanied by uncomfortable sensations, that begins or worsens during rest, is relieved by movement, and occurs predominantly in the evening or night. The condition causes significant distress or impairment. It is associated with iron deficiency, renal failure, pregnancy, and peripheral neuropathy. Ferritin levels below 50 ng/mL should be supplemented. Dopaminergic agents and alpha-2-delta ligands (gabapentin enacarbil) are used for treatment.

\medterm{Schizoaffective Disorder}-- A disorder characterized by an uninterrupted period of illness during which there is a major mood episode concurrent with criterion A of schizophrenia (delusions, hallucinations, disorganized speech, grossly disorganized or catatonic behavior, or negative symptoms). The individual must also have delusions or hallucinations for at least two weeks in the absence of a major mood episode. The mood episodes are present for the majority of the total duration of the illness.

\medterm{Schizoid Personality Disorder}-- A pattern of detachment from social relationships with a restricted range of emotional expression. At least four of seven criteria: neither desires nor enjoys close relationships, almost always chooses solitary activities, has little interest in sexual experiences, takes pleasure in few activities, lacks close friends, appears indifferent to praise or criticism, and shows emotional coldness, detachment, or flattened affect.

\medterm{Schizophrenia}-- A severe chronic mental illness characterized by psychotic symptoms (hallucinations, delusions, disorganized thinking) and negative symptoms (avolition, alogia, anhedonia, blunted affect), with cognitive impairment and functional decline. Diagnosis: \(\ge\)2 core symptoms for \(\ge\)1 month with continuous signs for \(\ge\)6 months. Onset typically in late adolescence/early adulthood. Treatment: antipsychotic medications (first-generation: haloperidol; second-generation: risperidone, olanzapine, clozapine for refractory), psychosocial interventions, and family education.

\medterm{Schizotypal Personality Disorder}-- A pattern of social and interpersonal deficits marked by acute discomfort with and reduced capacity for close relationships, cognitive or perceptual distortions, and eccentricities of behavior. Features include ideas of reference, odd beliefs or magical thinking, unusual perceptual experiences, odd thinking and speech, suspiciousness, inappropriate affect, and constricted affect. It shares genetic and phenomenological features with schizophrenia.

\medterm{Seasonal Affective Disorder}-- A specifier of recurrent major depressive episodes that occur at a specific time of year, most commonly winter (with remission in spring/summer). Symptoms: hypersomnia, hyperphagia with carbohydrate craving, weight gain, fatigue, and low energy. Associated with reduced sunlight exposure disrupting circadian rhythms and serotonin regulation. Treatment: light therapy (10,000 lux, 30 minutes upon waking), CBT-SAD, SSRIs/bupropion, and agomelatine. Dawn simulation and vitamin D supplementation may help.

\medterm{Sedatives}-- See Sedatives in Pharmacology.

\medterm{Selective Mutism}-- Consistent failure to speak in specific social situations where speaking is expected despite speaking in other situations. The disturbance interferes with educational or occupational achievement or social communication. Must last for at least one month and is not attributable to lack of knowledge of or comfort with the spoken language. Most commonly seen in children and often co-occurs with social anxiety disorder.

\medterm{Separation Anxiety Disorder}-- Developmentally inappropriate and excessive fear or anxiety concerning separation from attachment figures. The individual may have recurrent excessive distress when anticipating or experiencing separation from home or major attachment figures. Symptoms may include worry about harm befalling attachment figures, worry about experiencing an event causing separation, reluctance to leave home or go out alone, and nightmares involving separation. Must persist for four weeks in children and six months in adults. Diagnosis is clinical using DSM-5 criteria. Treatment includes CBT with graded exposure to separation situations and SSRIs (fluoxetine, sertraline, fluvoxamine) for moderate-to-severe cases. Involving parents in treatment is essential. Prognosis is good with treatment; untreated separation anxiety may persist and predispose to other anxiety disorders in adulthood.

\medterm{Serotonin-Norepinephrine-Dopamine Reuptake Inhibitors}-- A class of antidepressants that inhibit reuptake of serotonin, norepinephrine, and dopamine. The primary agent is bupropion, which also acts as a norepinephrine-dopamine reuptake inhibitor. It is used for depression, smoking cessation, and as augmentation. It lacks sexual side effects and weight gain associated with SSRIs.

\medterm{Serotonin-Norepinephrine Reuptake Inhibitors}-- Antidepressants that block the reuptake of both serotonin and norepinephrine. They include venlafaxine, duloxetine, desvenlafaxine, and levomilnacipran. SNRIs are used for major depressive disorder, generalized anxiety disorder, diabetic neuropathy, fibromyalgia, and chronic pain. They are particularly useful in patients with fatigue and poor concentration. Side effects include those of SSRIs plus hypertension (especially at higher doses).

\medterm{Sexual Masochism Disorder}-- A paraphilic disorder characterized by recurrent, intense sexual arousal from the act of being humiliated, beaten, bound, or otherwise made to suffer, lasting \(\ge\)6 months and causing distress or impairment or involving acting on these urges. Asphyxiophilia (sexual arousal from oxygen restriction): the most dangerous form, associated with accidental death during solo practice. Hypoxyphilia: arousal from oxygen deprivation. Consensual BDSM activities are not considered pathological. Disorder is diagnosed when the behavior causes clinically significant distress, impairment, or risk of harm. Treatment: CBT and SSRIs.

\medterm{Sexual Sadism Disorder}-- A paraphilic disorder characterized by recurrent, intense sexual arousal from the physical or psychological suffering of another person, lasting \(\ge\)6 months and causing distress or impairment or involving acting on these urges with a nonconsenting person. In the context of BDSM with consenting adults, sadistic acts alone do not constitute a disorder. Diagnosis requires that the individual has acted on these urges with a nonconsenting person or the urges cause significant distress or interpersonal difficulty. The most severe form is associated with serial sexual homicide. Treatment: CBT, SSRIs, antiandrogens, and GnRH agonists for severe cases.

\medterm{Social Anxiety Disorder}-- Marked fear or anxiety about one or more social situations in which the individual may be scrutinized by others. The individual fears acting in a way or showing anxiety symptoms that will be negatively evaluated. The social situations almost always provoke fear or anxiety, are avoided or endured with intense anxiety, and are disproportionate to the actual threat. The fear, anxiety, or avoidance is persistent, typically lasting six months or more.

\medterm{Somatic Symptom Disorder}-- Characterized by one or more somatic symptoms that are distressing or result in significant disruption of daily life, accompanied by excessive thoughts, feelings, or behaviors related to the somatic symptoms. The individual experiences persistent high levels of worry about the seriousness of symptoms, an excessive amount of time and energy devoted to health concerns, or disproportionate and persistent concern about the medical implications of symptoms. The problematic symptoms are not better explained by another medical or psychiatric condition. Somatic symptom disorder is common, with prevalence of approximately 5-7 percent in the general population. It is associated with high healthcare utilisation and functional impairment. Diagnosis is clinical using DSM-5 criteria, after appropriate medical evaluation to exclude underlying organic pathology. Treatment includes establishing a consistent therapeutic relationship with a single primary care provider, CBT, graded exercise, and antidepressants (SSRIs, TCAs) for comorbid depression or pain. Prognosis is variable, but most patients improve with appropriate management that avoids unnecessary investigations and procedures.

\medterm{Specific Learning Disorder}-- Persistent difficulties in reading, written expression, or mathematics, despite adequate intelligence and educational opportunity. Diagnosis requires academic skills significantly below expected for age, difficulties for at least six months, and no better explained by intellectual disability, sensory impairment, or neurological condition. Subtypes include dyslexia (reading), dysgraphia (writing), and dyscalculia (mathematics).

\medterm{Specific Phobia}-- Marked fear or anxiety about a specific object or situation that is out of proportion to the actual danger posed. Exposure to the phobic stimulus almost invariably provokes an immediate anxiety response. The object or situation is actively avoided or endured with intense anxiety. The fear, anxiety, or avoidance persists for six months or more. Subtypes include animal, natural environment, blood-injection-injury, situational, and other.

\medterm{Splitting}-- A primitive defense mechanism in which objects or self-representations are viewed as entirely good or entirely bad, with no integration of positive and negative qualities. It is a core feature of borderline personality disorder and contributes to idealization-devaluation patterns in relationships. It reflects impaired object relations.

\medterm{Stimulant Use Disorder}-- A problematic pattern of stimulant use (cocaine, amphetamines, methamphetamine, MDMA) leading to significant impairment. Stimulants increase dopamine and norepinephrine activity. Withdrawal is characterized by fatigue, depression, increased appetite, hypersomnia, psychomotor retardation, and vivid unpleasant dreams. Methamphetamine use is associated with severe dental disease, paranoia, and hallucinations. Treatment is primarily behavioral as no FDA-approved medications exist.

\medterm{Stress}-- A physiological and psychological response to perceived threats or demands that activates the sympathetic nervous system and the hypothalamic-pituitary-adrenal axis. Acute stress can enhance focus and performance in challenging situations, but chronic stress contributes to the development or exacerbation of numerous health conditions, including hypertension, anxiety disorders, depression, gastrointestinal problems, and immunosuppression. Effective stress management techniques include regular physical activity, mindfulness meditation, adequate sleep, social support, and professional counseling when needed.

\medterm{Substance Abuse Psychiatry}-- The subspecialty focused on the evaluation and treatment of substance use disorders and co-occurring psychiatric conditions. Dual diagnosis patients require integrated treatment addressing both disorders simultaneously. Psychopharmacology includes medication-assisted treatment (methadone, buprenorphine, naltrexone for opioid use disorder; acamprosate, disulfiram, naltrexone for alcohol use disorder) alongside psychotherapy including motivational interviewing and CBT.

\medterm{Substance/Medication-Induced Sexual Dysfunction}-- Sexual dysfunction directly caused by substance use or medication effects. Common culprits: SSRIs/SNRIs (decreased desire, delayed ejaculation/anorgasmia, erectile dysfunction — most common cause of drug-induced sexual dysfunction), antipsychotics (hyperprolactinemia causing decreased desire and ED), antihypertensives (beta-blockers, thiazides, alpha-blockers), hormonal contraceptives (decreased desire, vaginal dryness), antiandrogens (spironolactone, finasteride), opioids (decreased desire, ED, delayed ejaculation), and alcohol (acute: erectile difficulty; chronic: decreased desire and testicular atrophy). Management: identify the culprit, dose reduction, drug holiday, switch to alternative (bupropion, mirtazapine for depression; PDE5 inhibitors for ED), and add-on therapy (buspirone for SSRI-induced sexual dysfunction). Bupropion has the lowest rate of sexual side effects among antidepressants.

\medterm{Substance Use Disorder}-- A problematic pattern of substance use leading to clinically significant impairment or distress, manifested by at least two of eleven criteria within a twelve-month period. Criteria include taking the substance in larger amounts or over a longer period than intended, persistent desire or unsuccessful efforts to cut down, great deal of time spent obtaining or recovering from the substance, craving, failure to fulfill obligations, continued use despite social problems, giving up activities, recurrent use in hazardous situations, continued use despite physical or psychological problems, tolerance, and withdrawal. Severity is graded as mild (2-3 criteria), moderate (4-5), or severe (6 or more). Diagnosis is clinical using DSM-5 criteria. Treatment involves motivational interviewing to enhance motivation for change, CBT, contingency management, 12-step facilitation, and medication-assisted treatment where available (methadone, buprenorphine, naltrexone for opioid use disorder; acamprosate, disulfiram, naltrexone for alcohol use disorder). Relapse rates are high, and recovery often requires multiple treatment episodes and long-term support.

\medterm{Suicidal Ideation}-- Thoughts about killing oneself, ranging from passive (wish to be dead, wishes to fall asleep and not wake) to active (thoughts of killing oneself with varying degrees of intent, plan, and means). The C-SSRS grades severity from 1 (wish to be dead) to 5 (active suicidal ideation with specific plan and intent). Assessment of suicidal ideation must include questions about plan, intent, means, and history of attempts.

\medterm{Suicide}-- The intentional act of taking one's own life. Suicide is a major public health problem worldwide, causing ~700,000 deaths annually. Risk factors: psychiatric disorders (major depression—most common contributor, bipolar disorder, schizophrenia, borderline personality disorder, PTSD, substance use disorders—especially alcohol), previous suicide attempt (strongest predictor), family history of suicide, access to means (firearms—most common method in the US; hanging, pesticide poisoning—common globally), male sex (3--4:1 ratio of completed suicide vs attempts), older age (peaks in middle-aged and older adults), chronic physical illness, social isolation, unemployment, financial stress, history of childhood abuse/trauma, and LGBTQ+ identity (increased risk in the context of minority stress). Assessment: direct questioning about suicidal thoughts, plans, intent, and means (does not increase risk—contrary to myth). The SAFE-T (Suicide Assessment Five-step Evaluation and Triage) framework. Acute management: ensure safety—remove means, hospitalisation (voluntary or involuntary if imminent risk), suicide watch (1:1 observation), and pharmacological treatment of underlying psychiatric condition. Long-term management: psychotherapy (CBT, dialectical behaviour therapy—DBT for borderline personality disorder), pharmacotherapy (antidepressants—SSRIs, lithium—reduces suicide risk in mood disorders, clozapine—for schizophrenia), and psychosocial interventions (safety planning, crisis hotlines—988 in the US, 111 option 2 in the UK). Postvention: support for bereaved family and friends (suicide bereavement). Prevention: restricting access to means, responsible media reporting, school-based programmes, and gatekeeper training (GPs, teachers, police).

\medterm{Tardive Dyskinesia}-- Involuntary, repetitive, purposeless movements caused by prolonged use of antipsychotic medications, typically developing after months to years of exposure. Movements include choreoathetoid movements of the tongue, face, and extremities (lip smacking, tongue protrusion, grimacing, choreiform movements of limbs). It is caused by dopamine receptor supersensitivity. It may be irreversible. VMAT2 inhibitors (valbenazine, deutetrabenazine) are FDA-approved treatments. Risk is lower with atypical antipsychotics (especially clozapine and quetiapine) compared to typical antipsychotics. Prevention involves using the lowest effective dose, regular screening using the Abnormal Involuntary Movement Scale (AIMS) every 6-12 months, and switching to an atypical antipsychotic with lower D2 affinity if emerging dyskinesia is detected. Patients should be warned about the risk prior to starting antipsychotic treatment.

\medterm{Therapeutic Alliance}-- The collaborative, trusting relationship between therapist and patient that is foundational to effective psychotherapy. It includes agreement on goals, tasks, and the bond between therapist and patient. It is the strongest predictor of psychotherapy outcome across all therapeutic modalities.

\medterm{Therapeutic Drug Monitoring}-- The clinical practice of measuring drug levels in blood to ensure therapeutic concentrations while avoiding toxicity. It is particularly important for medications with narrow therapeutic windows including lithium (0.6--1.2 mEq/L), valproate (50--125 mcg/mL), carbamazepine (4--12 mcg/mL), and clozapine (350--600 ng/mL).

\medterm{Therapeutic Hypothermia}-- Controlled cooling of the body temperature used in psychiatric practice for severe catatonia and neuroleptic malignant syndrome. In NMS, cooling is a key treatment alongside discontinuation of the offending agent and dantrolene administration. Body temperature is reduced to below 38 degrees Celsius using cooling blankets, ice packs, or intravascular cooling catheters.

\medterm{Thought Blocking}-- An abrupt interruption in the train of thought before it is completed, as if the thought was removed or pushed out of the patient's mind. The patient may pause mid-sentence, appear confused, and report that the thought was ``taken out'' or ``blocked.'' It is a formal thought disorder seen in schizophrenia.

\medterm{Tourette Disorder}-- A neurodevelopmental disorder characterized by multiple motor tics and one or more vocal tics lasting >1 year, with onset before age 18. Motor tics: eye blinking, head jerking, shoulder shrugging, facial grimacing. Vocal tics: throat clearing, sniffing, grunting, barking, coprolalia (obscene words, occurs in <10\%). Premonitory urge precedes tics. Comorbidities: ADHD (60\%), OCD (30-50\%), anxiety. Course: waxing and waning severity; most improve by late adolescence. Diagnosis: clinical DSM-5 criteria. Treatment: CBIT (habit reversal training), alpha-2 agonists (clonidine, guanfacine), antipsychotics (haloperidol, pimozide) for severe tics.

\medterm{Tourette Syndrome}-- A neurodevelopmental disorder characterized by multiple motor tics and at least one vocal tic occurring over at least one year. Tics are sudden, rapid, recurrent, stereotyped movements or vocalizations. They are often preceded by premonitory urges. Tics typically worsen in adolescence and may improve in adulthood. Comorbidities include ADHD, OCD, and anxiety disorders. Treatments include behavioral therapy and antipsychotics for severe tics.

\medterm{Tranquillisers}-- See Tranquillisers in Pharmacology.

\medterm{Transference}-- A psychoanalytic concept in which feelings, attitudes, and expectations originally directed toward a significant figure from childhood are unconsciously redirected onto the therapist. In psychodynamic therapy, transference is a therapeutic tool for understanding the patient's relational patterns and resolving conflicts.

\medterm{Transvestic Disorder}-- A paraphilic disorder characterized by recurrent, intense sexual arousal from cross-dressing, lasting \(\ge\)6 months and causing marked distress or impairment. Almost exclusively in heterosexual males. Transvestic disorder is distinguished from gender dysphoria by the primary motivation: sexual arousal (transvestic) vs. identity incongruence (gender dysphoria). Many cross-dressers do not experience distress and therefore do not meet disorder criteria. Treatment: CBT for distress or guilt, and SSRIs for compulsive behavior. No treatment is indicated for cross-dressing without distress.

\medterm{Trichotillomania}-- Recurrent pulling out of one's hair, resulting in hair loss. The individual has made repeated attempts to decrease or stop hair pulling. The hair pulling causes clinically significant distress or impairment. It is not attributable to another medical condition. The pulling can occur from any body region and may be focused or focused-pulling. It is classified under obsessive-compulsive and related disorders.

\medterm{Tricyclic Antidepressants}-- Older antidepressants that non-selectively inhibit the reuptake of serotonin and norepinephrine. They include amitriptyline, nortriptyline, imipramine, clomipramine, and desipramine. TCAs are effective but have more side effects than SSRIs, including anticholinergic effects (dry mouth, urinary retention, constipation), orthostatic hypotension, weight gain, and cardiac conduction delays. Lethal in overdose due to cardiac arrhythmias. They are used for depression, neuropathic pain, migraine prophylaxis, and off-label for insomnia and anxiety. TCAs have significant anticholinergic, antihistaminergic, and alpha-adrenergic blocking properties. Nortriptyline has the least anticholinergic side effects and is preferred in elderly patients. Due to their narrow therapeutic index and cardiotoxicity in overdose, TCAs are used less frequently than SSRIs but remain important for treatment-resistant depression. Monitoring of ECG and serum drug levels is recommended.

\medterm{Vagus Nerve Stimulation}-- A neuromodulation treatment involving electrical stimulation of the left vagus nerve via an implanted device. It is FDA-approved for treatment-resistant depression and epilepsy. It requires surgical implantation and delivers intermittent stimulation. Improvement may take months to develop. Non-invasive transcutaneous vagus nerve stimulation is being investigated.

\medterm{Voyeuristic Disorder}-- A paraphilic disorder involving recurrent, intense sexual arousal from observing an unsuspecting person who is naked, disrobing, or engaging in sexual activity. The individual has acted on these urges or they cause marked distress or interpersonal difficulty. Duration \(\ge\)6 months. Most common in males with onset in adolescence. Treatment: CBT and SSRIs.

\medterm{Word Salad}-- Severe disorganized speech consisting of a jumble of words and phrases that lack meaning or logical connection. It is a hallmark of formal thought disorder in schizophrenia and represents the most severe form of disorganized speech. It reflects profound disruption of language processing.

